\begin{statementbox}
\textbf{C012\quad 椭圆切线原心距性质}

设椭圆
\[
C:\quad \frac{x^2}{a^2}+\frac{y^2}{b^2}=1,
\qquad a>b>0,\quad c^2=a^2-b^2,\quad e=\frac ca.
\]
其中心为 $O(0,0)$, 焦点为 $F_1(-c,0),F_2(c,0)$.
任取椭圆上一点 $P(x_0,y_0)$, 在 $P$ 处作切线 $l$.
记 $r_1=PF_1$, $r_2=PF_2$, $d=\operatorname{dist}(O,l)$.
则
\[
\boxed{\sqrt{r_1r_2}\,d=ab}
\qquad\Longleftrightarrow\qquad
\boxed{d=\frac{ab}{\sqrt{r_1r_2}}}.
\]
即\textbf{切点到两焦点距离的几何平均数, 乘以中心到切线的垂直距离, 恒等于 $ab$}.

\textbf{常用等价式:}
\[
r_1r_2=a^2-e^2x_0^2,
\qquad
\boxed{d^2=\frac{a^2b^2}{a^2-e^2x_0^2}}.
\]
由于 $b^2\le r_1r_2\le a^2$, 必有
\[
\boxed{b\le d\le a}.
\]
其中 $d=a$ 当且仅当 $P=(\pm a,0)$ (长轴顶点);
$d=b$ 当且仅当 $P=(0,\pm b)$ (短轴顶点).

\textbf{几何补充:} 若 $\varphi=\angle F_1PF_2$, 则
\[
\boxed{d=a\cos\frac\varphi2},\qquad
\boxed{r_1r_2\cos^2\frac\varphi2=b^2}.
\]
上述两式也是本结论的纯几何表达.
\end{statementbox}
